 %!TEX root = main.tex

 \subsection{Response Strategy Framework}

For each adverse scenario, three response strategies are evaluated:

\begin{itemize}[leftmargin=*,itemsep=2pt]
  \item \textbf{R : Rerouting only.} The baseline infrastructure
    ($W^*, A^*$) is held fixed; only flow routing is re-optimised.
    Fixed costs are sunk and not re-incurred. This represents the
    immediate, operational response horizon.

  \item \textbf{A : Adaptation.} Warehouse opening and arc activation
    decisions are freed; the solver may open or close facilities and
    activate or deactivate arcs at their respective costs. This represents
    a medium-term reconfiguration.

  \item \textbf{F : Full redesign.} Phase~1 is re-solved from scratch under
    the scenario's parameters, ignoring sunk costs. This gives a lower bound
    on achievable cost and represents the long-run strategic horizon.
\end{itemize}

The disruption cost for a given scenario~$s$ and strategy~$\sigma$ is:
\begin{equation}
  \Delta Z(s,\sigma) = Z(s,\sigma) - Z^*
  \label{eq:deltaz}
\end{equation}
where $Z^*$ is the baseline optimal cost from Phase~1.

\subsection{Scenario Overview}

Table~\ref{tab:scenarios} summarises the six scenarios.

\begin{table}[H]
\centering
\caption{Phase~2 adversarial scenarios.}
\label{tab:scenarios}
\begin{tabular}{lll}
\toprule
\textbf{ID} & \textbf{Type} & \textbf{Shock} \\
\midrule
T1 & Tariff      & $+40\%$ transatlantic/transpacific \\
T2 & Energy      & Air $+60\%$, road $+20\%$, sea $+10\%$ \\
T3 & Seasonal    & Hub handling fees $+30\%$ (Q4) \\
S1 & Node loss   & H3 (Singapore) closed \\
S2 & Arc loss    & Korea--Europe corridor blocked \\
S3 & Combined    & T2 $+$ S1 simultaneously \\
\bottomrule
\end{tabular}
\end{table}

\paragraph{T1 : Trump-era tariff shock.}
A 40\% surcharge is applied to all arcs crossing the transatlantic and
transpacific corridors, primarily affecting products~B and~C whose supply
chains rely heavily on Asia--Americas and Asia--Europe routes.
Under strategy~R, flows can be partially rerouted through Middle East or
intra-regional corridors; under~F, the network may reorganise around
lower-tariff paths.

\paragraph{T2 : Energy crisis.}
Inspired by the Ukraine/Iran conflict, this scenario raises air freight by
60\%, road by 20\%, and sea by 10\%. The asymmetric impact penalises
time-sensitive, air-intensive routing heavily. Products with modal
flexibility benefit most from adaptation.

\paragraph{T3 : Seasonal hub surcharge.}
A 30\% surcharge on hub handling fees during Q4 raises the effective cost
of all hub-transiting flows. Products that route through multiple hubs
(notably~A via sea and~B via air) are most exposed.

\paragraph{S2 : Korea--Europe corridor blockade.}
A set of arcs representing the South Korea export corridor to Europe are
removed. This primarily disrupts product~B (semiconductors), which relies
on Korean-origin flows. Alternative routing through other hubs is feasible,
at higher cost.

\subsection{S1 and S3: Structural Infeasibility and Modified Formulation}
\label{sec:s1s3}

Scenarios S1 and~S3 involve the closure of hub~H3 (Singapore), motivated
by a simulated Taiwan Strait crisis. Unlike the other scenarios, these
are \emph{structurally infeasible} under the standard formulation
(equations \ref{eq:C1}--\ref{eq:C8}).

\paragraph{Root cause.}
Suppliers S4, S5, and~S6 (producing product~B) and supplier~S9
(producing product~C) are exclusively connected to the network through~H3.
When H3 is removed, all 42 arcs incident to it are severed, leaving
these suppliers with no outgoing path. The resulting supply deficit is:

\begin{itemize}[leftmargin=*,itemsep=1pt]
  \item Product~B: $-2{,}223$ units (100\% of total supply stranded)
  \item Product~C: $-631$ units (33\% of total supply stranded)
\end{itemize}

With C1 as a hard equality constraint, no feasible solution exists.

\paragraph{Modified formulation for strategies R and A.}
We introduce a non-negative slack variable $unmet_c^p \in [0, Dem_c^p]$
representing unserved demand, and replace~\eqref{eq:C1} with:
\begin{equation}
  \sum_{a \in S_{\text{target}}^c} x_a^p + unmet_c^p = Dem_c^p
  \quad \forall\, p \in P,\; c \in C
  \label{eq:C1_relaxed}
\end{equation}
The objective gains a penalty term $M \cdot \sum_{c,p} unmet_c^p$ with
$M = 10^6$, large enough to ensure unmet demand is never preferred over
any feasible routing. The logistics cost reported excludes the penalty:
$Z_{\text{logistics}} = Z_{\text{obj}} - M \cdot \sum_{c,p} unmet_c^p$.

\paragraph{Modified formulation for strategy F.}
Full demand satisfaction
is restored by injecting one emergency supply arc per
stranded supplier, connecting it to the geographically
closest remaining active hub. The unit cost scales with
distance: $c_a^{\text{emg},p} = \gamma \cdot \hat{c}^{s,p}
\cdot \hat{d}^{s,h}$, where $\hat{c}^{s,p}$ is the
cost-per-km derived from the supplier's original baseline
arc, $\hat{d}^{s,h}$ is the estimated rerouting distance
via the removed hub, and $\gamma = 2.5$ models spot-market
procurement. Under~F, constraint~\eqref{eq:C1} is enforced
as a hard equality~--- the emergency arcs restore
feasibility while the objective captures the full
procurement premium.

